% !TEX root = Master Thesis BWL.tex
%\pagestyle{fancy} % diese und n�chsten 4 Zeilen f�r sch�ne Kopf- und Fu�zeile
%\fancyhead[LO,ER]{Introduction} % [left-uneven,right-even] {Chaptername}
%\fancyhead[RO,EL]{\thepage} % [right-uneven, left-even]{pagenumber}
%\renewcommand{\headrulewidth}{0.4pt} % unterhalb der Kopfzeile kommt ein Trennstrich

%\newgeometry{top = 2.54 cm, bottom = 2.54 cm}
%\footskip = 35pt

\chapter{Introduction}
Cycling races are among the most exciting and popular events in the sports world, attracting millions of fans and participants worldwide. From the Tour de France to the Giro d'Italia, these races are well known by vivid cycling advocates and even by non-cyclists. Those races offer a unique combination of strategy, endurance, and teamwork, making them a fascinating subject of study for game theorists.
Game theory is a branch of mathematics that deals with the strategic interactions between decision-makers, where each player's outcome depends not only on their actions but also on the actions of others. It has been applied to various fields, including economics, political science, and biology, and has proven to be a powerful tool for analyzing complex systems and designing optimal strategies \cite{MyersonRogerB2004Gt:a}. In the context of cycling races, game theory provides a framework for analyzing the strategic interactions between the riders and the teams. Each rider and team is a decision-maker who must choose their strategies based on the actions of others. The goal is to identify the optimal strategies that enable them to achieve their objectives in the race, such as winning the race.
\newline
In this Master's thesis, we begin our analysis by focusing on the smallest game-theoretically relevant unit of a cycling race and gradually build up to the individual races of a complex Omnium event. Cycling races are complex events that involve not only physical abilities but also strategic decision-making. Game theory provides a powerful framework to study the strategic interactions among cyclists and the various factors that affect their performance. By starting with the smallest unit and gradually increasing the complexity, the aim is to gain a comprehensive understanding of the strategic dynamics of cycling races and their implications for race outcomes.
\newline
Overall, this thesis aims to contribute to the understanding of the strategic interactions in cycling races. By applying game theory to the world of cycling, we hope to shed light on the complex dynamics of this exciting sport and provide a new perspective on the tactics and strategies used by the riders.
\section{Structure of the Thesis}
To provide a clearer understanding of the thesis, firstly the specific game-theoretical aspects relevant to cycling will be explored. The first section gives an overview of the cycling world, examining the various types of riders within a cycling team and delving into the strategic considerations behind selecting certain types based on the team's goals for the upcoming cycling season. Additionally, the discussion will encompass an analysis of different event types, specifically those that will be examined in detail within the thesis.The focal point of the thesis lies in the second chapter, titled ``Strategic Behavior in Cycling.'' In this chapter, the focus will shift towards exploring game theoretical models that can be applied to comprehend a range of scenarios, from analyzing individual attacks, which are fundamental in cycling, to understanding entire events like the Omnium. However, it is important to note that the thesis will not delve into the intricacies of the Omnium itself, which can be considered a game within a game. This is due to the fact that such a detailed analysis would surpass the scope of this study.
\section{Motivation and personal involvement}
Blending game theory principles with sports isn't something entirely new, especially in a field often referred to as ``chess on wheels'' - road cycling. As someone deeply entrenched in the world of competitive cycling at a national level, I grasp the intricate strategies that fuel this sport. It's more than just turning pedals; it's a domain rich in strategic thinking, clever tactics, and the art of seamless teamwork. In this vibrant sport, the ability to anticipate your opponent's next move is often the difference between victory and being left behind in the peloton. Cycling demands mental strength as much as physical endurance. Over years of trials and errors in the saddle, cyclists acquire the invaluable skill of tactical thinking. This thesis aims to shed light on the strategic nuances that define competitive cycling.
A major source of inspiration for my master's thesis comes from Jean-Francois Mignot's 2015 paper, ``Strategic Behavior in Road Cycling Competitions.'' This paper, along with insights from my personal racing experiences and a close collaboration with a professional German cycling team, has given me the confidence to dive into this topic.
My knowledge didn't come solely from personal endeavors; it was significantly enhanced through discussions with experienced professionals, each boasting a decade or more of expertise in the field. These conversations were important, providing me with knowledge that directly influenced this research. Combining my own experiences on the race track, insights shared by professionals, and rigorous research, this thesis aims to provide a comprehensive overview of the strategic tactics prevalent in road cycling.
\section{Cycling}
Cycling is a sport and recreational activity that involves riding a bicycle on roads, tracks, or off-road trails. It can be done for leisure, fitness, transportation, or competition. Cyclists use a variety of bicycles, including road bikes, mountain bikes, and hybrid bikes, to travel long or short distances, depending on their goals and preferences.
In competitive cycling, riders compete in races that vary in length, terrain, and difficulty. Road cycling is one of the most popular forms of competitive cycling, which involves racing on paved roads, often over long distances. Other forms of competitive cycling include track cycling, cyclocross, mountain biking, and BMX racing.
The thesis concentrates purely on road - and track cycling. The next sections of this master thesis provide a comprehensive analysis of road cycling from various angles, exploring the different types of cycling events. Additionally, we delve into the types of riders existing in a competitive field. Finally, we examine the science of drafting, exploring the benefits to gain a competitive edge. Through this analysis, we hope to gain a deeper understanding of the complexities and nuances of road cycling to have a base knowledge to understand the strategies and tactics that riders and teams use to succeed in this highly competitive sport \cite{mignot2015history}.

\subsection{Event types}
In general, cycling is divided into two types of events. One-day races, which take place on one day only, without a race the day before or after, and multi-stage races, which take place at least two days in a row. Both kinds present different challenges that must be overcome by the riders. This chapter will discuss the two types and their challenges.
\subsubsection{Multi-Stage Races}
Multistage races are a type of cycling race that typically takes place over several days and covers hundreds or even thousands of kilometers.
They are some of the most grueling and demanding events in the sport, requiring riders to have exceptional endurance,
tactical awareness, and mental strength. One of the most prestigious and well-known multistage races is the Tour de France,
which takes place over three weeks in July and covers over 3,000 kilometers.
The race is divided into 21 stages, each with its own unique characteristics, such as mountain climbs,
flat sprints, or time trials. Other notable multistage races include the Giro d'Italia and the Vuelta a Espa\~na,
both of which are held in Europe and cover similar distances over a period of several weeks.
\newline
In multistage races, riders compete as individuals and as part of a team. Each stage is a separate race, with the rider or team that completes the course in the shortest time declared the winner. Points are also awarded for finishing positions in each stage, and the rider with the most points at the end of the race is awarded the overall winner's jersey. A team usually has a rider that is designated to win the general classification (the whole stage race) and another rider going for the stage win. Riders that are picked by their ``directors sportive'' to win, are to be protected at all costs. They are picked by looking at previous race, results and general fitness at the moment of the race. The rest of the team has the task of bringing the designated winner to the finish line as fresh as possible.
Discussing the strategic intricacies of multi-stage races in the thesis would exceed its intended scope. Each stage of the race has its unique characteristics and requires examination as both an individual component and as part of the entire event. The complexity arises from the multitude of details that demand attention. We believe it is more crucial to focus on laying the groundwork rather than delving into every nuanced aspect.
\subsubsection{One-Day Events}
Looking at the world tour one-day Events are usually reffered to as the ``Classics.'' These races are typically held on a single day and cover a distance of  200-300 kilometers. One-day races are known for their challenging terrain, with steep climbs and technical descents often included in the route. Popular one-day races include the Tour of Flanders, Paris-Roubaix, and the Milan-San Remo.
For the category elite and continental one-day events are very popular as most cyclists have to work besides racing their bikes. One-day races can be divided into different types or categories. One category is the circuit race, which involves a circular course that must be completed multiple times. Another one would be point-to-point races where start and finish are at different locations.
\newline
To win a one-day race, teams must employ a range of tactics that vary depending on the race's specific characteristics. On hilly terrain for example, climbers (usually small and light cyclists) may be sent to the front to attack and thin out the field. On flat terrain, teams may aim to control the peloton's pace and set up a sprint finish for their designated sprinter. Tactics can also include setting up breakaways, where a small group of riders escapes the peloton and works together to maintain their lead.
\subsubsection{Track cycling}
Track cycling is a form of cycling that takes place on a track or velodrome, and can be held indoors or outdoors. Velodromes can vary in size and shape, with some presenting shorter, tighter curves and others featuring longer, more gradual curves. They can also vary in the number of lanes and the distance of the track, with tracks measuring anything between 130 and 500 meters with tracks measuring anything between 130 and 500 meters. Special bikes, known as track bikes or fixed-gear bikes, are used on the velodrome for track cycling. These bikes have a single gear and no brakes, which means that the rider must continually pedal to keep the bike moving. Additionally, they must rely on their legs and the resistance of the track to slow down and stop. The frame geometry and handlebar design of track bikes are unique allowing for optimal aerodynamics and handling on the velodrome. This category involves several different disciplines, including sprints, endurance races, and team events, of which the sprint is one of the most well-known. Here two riders compete head-to-head over a short distance of 200-250 meters. The sprint is all about explosive power and speed, with riders jockeying for position and trying to outmaneuver each other as they approach the finish line. Endurance races, on the other hand, can range from the relatively short 4-kilometer individual pursuit to the grueling 40-kilometer points race. Endurance races require a combination of speed, endurance, and tactical ability, with riders often working together in teams to strategize and outpace their competitors. In this thesis, we focus on two events - Omnium and Sprint - which are particularly interesting from a game theoretical perspective \cite{Trackcycling}.
\begin{figure}[h]
\centering
\includegraphics[width=8cm]{figures/velodrome.png}
\caption[Velodrome]{Velodrome \cite{BBCNews_2006}}
\centering
\end{figure}
\subsubsubsection{Omnium}
\label{chap:Omnium}
The Omnium is a multi-event discipline in track cycling that consists of several races held on a single day. It tests an athlete's overall ability, as it includes both endurance and sprinting races. The Omnium typically consists of four to six events. The order and specific events may vary. This thesis concentrates on four events within the Omnium discipline: the Points Race, Elimination Race, Tempo Race, and Scratch Race. The Points Race is a longer endurance event where riders accumulate points by placing in intermediate sprints and finishing positions. The Elimination Race is a high-speed race where the last rider to cross the finish line every few laps is eliminated until only one rider remains. The Tempo Race is a newer event where riders race for a set number of laps before a bell rings, signaling a sprint for double points. The Scratch Race is a straightforward race where the first rider across the finish line wins. The winner of each event is awarded a certain number of points, and at the end of the competition, the athlete with the most points is declared the winner of the Omnium. This format rewards consistency and versatility, as athletes must perform well across a range of different events to win \cite{dwyer2013elimination}.
\subsection{Teams and Riders}
Often, cycling is mistakenly considered an individual sport, but upon closer examination of its structures, it becomes evident that cycling is, in fact, a team sport. Various riders, each with their unique profiles, are assigned specific tasks to achieve the best possible outcome for the team. This concept is known as Team Reasoning. Within a cycling team, individual riders can identify race strategies that might not directly benefit them but lead to the best possible payoff for the team as a whole. In such instances, the individual rider is willing to sacrifice personal gains to contribute to the team's success. This selfless approach exemplifies the collaborative nature of cycling and underscores the importance of teamwork in achieving overall victory
 \cite{sugden1993thinking, sindik2008application}. Each rider has their unique strengths and specialties, and it's essential for the team to have a balanced palette of riders to increase the potential of winning races. The following are the seven types of riders typically found in a cycling team:
\begin{itemize}
\item Rouleur (All-Rounder): A rouleur is a versatile rider who can perform well in a variety of terrain, such as flat roads, rolling hills, and cobbled sections. They are often the workhorses of the team and are responsible for setting the pace in the peloton, chasing down breakaways, and protecting the team's leaders.
\item Grimpeur (Climber): Grimpeurs are riders who excel in climbing mountains and hills. They are lightweight and have a high power-to-weight ratio, allowing them to maintain a high pace even when the road gets steep. They are crucial in mountain stages of races and are often the team's leaders.

\item Puncheur (Puncher): Puncheurs are riders who can handle short, steep climbs and are powerful on hilly terrain. They have explosive power and can accelerate quickly, making them ideal for breakaways and attacking the peloton.

\item Sprinters: Sprinters are riders who specialize in fast, flat finishes. They have a high top speed and are often the team's designated finisher in mass sprint finishes.

\item GC Riders: GC (General Classification) riders are the team's leaders and are responsible for winning the overall race. They have a well-rounded skill set and can perform well in all types of terrain. They are often the strongest climbers and time trialists on the team.

\item Time Trialists: Time trialists excel in individual time trials, where riders race against the clock. They have a high power output and aerodynamic position, allowing them to maintain a high speed for extended periods.

\item Domestiques (Servants): Domestiques are riders who sacrifice their chances of winning to support their team leaders. They perform tasks such as fetching water bottles, protecting their leaders from the wind, and setting the pace in the peloton. They play a crucial role in the success of the team and are often unsung heroes.
\end{itemize}
Overall, having a balanced palette of riders is crucial for a team's success in cycling. Each rider plays a specific role in achieving the team's goal, whether it be winning a race or supporting their leader \cite{Ridertype}.In a paper written by Sindik and Vidak in 2008, the authors delve into the concept of payoff with predictability among opposing players and players within the same team. Now, let's apply this notion to the context of cycling. Undoubtedly, predictability concerning the actions of both rival teams and one's teammates plays a vital role in shaping strategic decisions in a race. Understanding the predictability of rival teams allows cyclists to anticipate their moves, enabling more effective counter-strategies. On the other hand, having insight into the behavior of their own teammates helps in coordinating efforts and optimizing team tactics to achieve the best possible outcomes. In the dynamic world of cycling, where split-second decisions can determine the race's outcome, the level of predictability becomes a key factor in strategy formation. By taking cues from the concepts explored in paper by Sindik and Vidak \cite{sindik2008application}, riders and teams can enhance their chances of success by making well-informed and calculated choices based on their understanding of the predictability aspects in the race environment.
\begin{table}
\centering
\setlength{\extrarowheight}{2pt}
\begin{tabular}{cc|c|c|}
  & \multicolumn{1}{c}{} & \multicolumn{2}{c}{Team $2$} \\
  & \multicolumn{1}{c}{} & \multicolumn{1}{c}{Team Rider}  & \multicolumn{1}{c}{Opponent Rider}   \\\cline{3-4}
Individual Rider   & Predictable & $10$ & $0$ \\\cline{3-4}
            & Unpredictable & $-10$ & $-10$ \\\cline{3-4}
\end{tabular}
\caption{Outcome Matrix predictability according to Jo�ko Sindik and Nives Vidak (2008)
 }
\label{Matrix predictability}
\end{table}
Being predictable to teammates but unpredictable to opponents can bring the best utility, as opponent riders are left confused about the tactics of the other team. However, in real life, this scenario is very unlikely as it is hard to predict what tactics might be unforeseeable to other teams. Being unpredictable to both co-riders and opponents is very unfavorable. This might lead to confusion within the team, which is the last thing a team wants during a race. Additionally, opponents would be left in confusion, which could be favorable for the team. The worst scenario is to be unpredictable for one's teammates but predictable for the opponents. Opponents could use this knowledge as an advantage to outplay the team. Being predictable for both teammates and opponents would lead to teammates having a potentially very good action plan for the race. Although opponent teams can plan the perfect countertactic, the countertactic cannot surprise the other team, hence its utility is 0 in the matrix. This is why it is generally better to be predictable for teammates and opponents. Being unpredictable and causing confusion in a peloton might also lead to more crashes and unsafe situations during races.
\subsection{Drafting}
Drafting is a critical aspect of cycling as it plays a vital role in conserving energy and reducing air resistance, ultimately impacting the outcome of a race. The positioning of a rider in a peloton or a breakaway group is a key determinant of their drafting effectiveness. The riders at the front of the group face the maximum wind resistance and must expend more energy to maintain their speed, whereas the riders behind them can reduce their effort by a lot due to the slipstream effect. In \ref{Drag in a Peloton} one can see that riders in the front of the peloton use more energy than riders in the middle to back of the peloton. The graphic showed that riders use only as little as 10 percent of energy compared to if they ride alone \cite{gibertini2008cycling}. Hence, maintaining a good position in the peloton, such as in the middle or near the front, is crucial in conserving energy and increasing drafting efficiency. However, as each rider has a unique body type, drafting preferences also vary, and riders must learn through years of practice to position themselves optimally in different scenarios. Therefore, the competition for good drafting positions is fierce, and gaining expertise in slipstreaming can provide a competitive advantage to a cyclist in a race. From a game theory perspective, drafting plays a critical role in strategic decision-making. For example, in a race where a breakaway group of riders has formed, the riders in the chasing group may need to work together to draft and take turns at the front of the group to catch up to the breakaway. This involves making strategic decisions about when to take turns, how much effort to exert, and how to communicate with one another. On the other hand, if a rider or team has a significant lead, they may strategically use drafting to conserve energy and maintain their advantage. In this case, the other riders or teams must decide whether to continue to work together to catch up or to conserve their energy and wait for a better opportunity to attack \cite{gibertini2008cycling}.
\begin{figure}[h]
\centering
\includegraphics[width=10cm]{figures/drag foto.jpg}
\caption[Drag]{Drag in a Peloton\cite{BLOCKEN2018319}}
\label{Drag in a Peloton}
\centering
\end{figure}s
